\chapter{讨论}
\label{ch:discussion}

\section{实验平台的功能正确性}

完整插件是显示优化实验的平台，不是单独的研究对象。它能够保存和运行行为树，也支持 Actor、复合节点、条件、动作、Decorator、黑板和实时调试。因此，实验使用的是真正能够运行的行为树，而不是只画在画布上的静态节点图。

资源与运行时、编辑器界面、基础游戏和复杂竞技场共 450 项检查全部通过。241 节点行为树能够控制敌人的索敌、防御、近战、远程攻击、跳跃、攀爬、搜索、撤退和治疗。这些结果说明插件能够为显示优化实验提供功能正确、具有实际游戏行为的测试平台。

\section{显示优化效果评估}

\subsection{不同节点规模下的优化效果}

Compact 在五种规模下保留全部卡片，同时把卡片总面积减少 61.09--61.35\%。各规模的降低比例非常接近，说明紧凑卡片的效果不会随着节点数量增加而明显减弱。它适合在仍需查看全部节点时减少单张卡片占用的空间。

Optimized Overview 同样保留全部卡片，面积降低稳定在 90.27--90.34\%。与 Compact 相比，它进一步减少卡片内部信息，因此更适合查看大型行为树的整体分支结构。31 到 364 节点之间的结果变化很小，说明该组合在不同树规模下都能保持相近的信息密度改善。

Optimized Search 的非目标淡化比例从 31 节点树的 96.15\% 增加到 364 节点树的 99.67\%。随着行为树变大，画面中的非目标节点更多，搜索对目标显著性的提升也更明显。实验中的固定目标在所有规模下都能进入视口，说明 Search 能够在大型树中稳定完成已知节点定位。

Subtree Focus 在五种规模下减少了 76.92--91.58\% 的显示卡片，364 节点树为 85.90\%。该比例由目标分支本身的大小决定，所以不会随总节点数严格递增。Context Collapse 的卡片减少比例从 69.23\% 增加到 96.39\%，说明树越大，可折叠的非目标分支越多。两种方法都能减少当前画面中的无关上下文，但 Focus 面向指定分支，Collapse 则保留更多整体结构位置。

综合五种规模的数据，Compact 和 Overview 的面积降低保持稳定，Search 的目标突出效果随规模增加，Focus 和 Collapse 则直接减少同时显示的上下文。这些结果说明显示优化在小型树中已经有效，在 121、241 和 364 节点的大型树中作用更加明显。

\subsection{不同物理屏幕尺寸下的优化效果}

屏幕尺寸实验显示，物理显示面积会直接影响大型树能够呈现的上下文。屏幕从 15.94 英寸增大到 31.55 英寸后，241 节点树的 Baseline 覆盖率从 18.32\% 增加到 46.53\%，364 节点树从 12.13\% 增加到 30.82\%。两种复杂树的覆盖率都变为小屏幕的 2.54 倍，说明更大的物理屏幕能够同时容纳更多节点。

不同功能对屏幕尺寸的反应并不相同。Compact 在三块屏幕上的卡片面积降低均为 61.09--61.35\%，Overview 均为 90.27--90.34\%。这两种功能在小屏幕和大屏幕上的相对降低比例几乎相同，因此它们提供的是不依赖屏幕尺寸的信息密度改进，并没有额外放大大屏幕的优势。在 241 和 364 节点树上，完整画布已经达到 0.10 倍的最小缩放，所以 Compact 和 Overview 虽然明显缩小卡片及其信息量，却没有让小屏幕一次显示更多完整树节点。

Subtree Focus 更明显地补偿了小屏幕的限制。241 节点树在小屏幕上的覆盖率由 Baseline 的 18.32\% 提高到聚焦分支的 82.35\%，增加 64.04 个百分点；在大屏幕上则由 46.53\% 提高到 88.24\%，增加 41.70 个百分点。小屏幕与大屏幕之间原有的 28.22 个百分点差距因此缩小到 5.88 个百分点。364 节点树也出现相同方向：小屏幕提高 73.92 个百分点，大屏幕提高 64.53 个百分点，屏幕间差距由 18.69 个百分点缩小到 9.30 个百分点。这里比较的是当前保留分支的覆盖率，不表示完整树已经全部显示。

Context Collapse 也能减少小屏幕上同时显示的无关分支，但补偿效果弱于 Focus。241 节点树的小屏幕与大屏幕差距由 28.22 个百分点缩小到 14.29 个百分点；364 节点树的差距则只由 18.69 缩小到 18.18 个百分点。这是因为 Collapse 仍保留折叠根和路径上下文，而 Focus 更彻底地限制当前显示范围。因此，当开发者只需要处理一个局部分支时，Focus 对小屏幕最有价值；需要保留整体位置关系时，Collapse 提供了较温和的折中。

Search 在全部 15 个“屏幕--规模”组合中都把固定目标移入视口，说明目标定位能力没有因屏幕变小而失效。该结果能够证明搜索在三种屏幕上都可靠，但没有记录真人寻找节点所需的时间，因此不能进一步声称小屏幕用户节省了多少操作时间。

综合来看，大屏幕仍然具有绝对优势：即使使用同一种优化，它通常仍能显示更多上下文。但是，优化没有把这种优势进一步扩大。Compact 和 Overview 的相对效果基本不受屏幕尺寸影响，Search 保持跨屏幕稳定，而 Focus 明显缩小了小屏幕与大屏幕的覆盖差距。因此，本实验更支持“显示优化能够补偿小屏幕的部分空间限制，同时保留大屏幕的绝对容量优势”，而不是“显示优化主要扩大大屏幕优势”。这一结论只针对卡片面积、上下文覆盖和目标定位等可测界面属性，不直接等同于人的可读性或使用效率。
